\section{稀溶液的依数性}


\begin{frame}{稀溶液的依数性}
    由不挥发溶质构成的溶液，有以下4个性质：
    \begin{itemize}
        \item[1.] 蒸气压下降
        \item[2.] 沸点升高
        \item[3.] 凝固点降低
        \item[4.] 在半透膜两侧能产生渗透压
    \end{itemize}
    如果是非电解质的稀溶液，这四种性质在数量上又仅与溶质的质点数有关，而与溶质的性质无关，这种性质称为稀溶液的依数性(coligative property)。
    \\稀溶液的蒸气压下降是拉乌尔定律的直接结果，由化学势可导出其他三种性质与稀溶液浓度间的关系。
\end{frame}


\begin{frame}{渗透压}
    在如图所示的装置中，用半透膜将纯溶剂与溶液隔开，分为I、II两部分，半透膜只允许溶剂分子穿过，而溶质分子不能穿过。在一定温度下，溶剂分子穿过半透膜进入溶液，这种现象称为渗透。当渗透达平衡时，溶液的液面会升高，由于液面升高而产生的额外压力$\pi$称为渗透压(osmotic pressure)。
    \\图（a）为渗透开始时的状态，（b）为渗透平衡时的状态。
    $$\includegraphics[height=3.6cm]{figs/fig2.3.png}$$
\end{frame}


\begin{frame}{渗透压}
    渗透开始，纯溶剂和溶液中溶剂A的化学势分别为
    $$\begin{array}{c}
        \mu_{\mathrm{A}}^{\mathrm{I}}(T, p)=\mu_{\mathrm{A}}^{*}(T, p) \\
        \mu_{\mathrm{A}}^{\mathrm{II}}(T, p)=\mu_{\mathrm{A}}^{*}(T, p)+R T \ln x_{\mathrm{A}}
    \end{array}$$
    因为$x_\mathrm{A}<1$，所以$\mu_{\mathrm{A}}^{\mathrm{I}}(T, p)>\mu_{\mathrm{A}}^{\mathrm{II}}(T, p)$，纯溶剂要向溶液方向扩散，产生渗透现象。当渗透达平衡时，溶液与纯溶剂的液面差引起的压力即为渗透压$\pi$。此时溶液上方的压力为$p+\pi$。因已达平衡，纯溶剂与溶液中溶剂的化学势相等。
    $$\mu_{\mathrm{A}}^{\mathrm{I}}(T, p)=\mu_{\mathrm{A}}^{\mathrm{II}}(T, p+\pi)$$
    经推导得
    $$\pi = cRT$$
\end{frame}


\begin{frame}{凝固点降低}
    固体溶剂与溶液成平衡时的温度称为溶液的凝固点(freezing point)，这里指由纯溶剂形成的固体，而不是固溶体。设在压力$p$时，溶液的凝固点为$T$，则在两相平衡时有
    $$\mu_{\mathrm{A}}(\mathrm{l}, T, p, x_{\mathrm{A}})=\mu_{\mathrm{A}}(\mathrm{s}, T, p)$$
    式中：l和s分别表示液态和固态。在定压下若使溶液的浓度由$x_{\mathrm{A}}$变到$x_{\mathrm{A}}+\mathrm{d}x_{\mathrm{A}}$，则凝固点相应地由$T$变到$T+\mathrm{d}T$，重新建立平衡。
    \\稀溶液的凝固点降低公式
    $$\Delta T_{\mathrm{f}}=\frac{R\left(T_{\mathrm{f}}^{*}\right)^{2}}{\Delta_{\mathrm{fus}} H_{\mathrm{m}}(\mathrm{A})} \frac{n_{\mathrm{B}}}{n_{\mathrm{A}}}$$
\end{frame}


\begin{frame}{凝固点降低}
    
    设溶剂和溶质的质量分别为$m_\mathrm{A}$和$m_\mathrm{B}$，其摩尔质量分别为$M_\mathrm{A}$和$M_\mathrm{B}$，则上式又可写成
    $$\Delta T_{\mathrm{f}}=\frac{R\left(T_{\mathrm{f}}^{*}\right)^{2}}{\Delta_{\mathrm{fus}} H_{\mathrm{m}}(\mathrm{A})} \frac{M_\mathrm{A}m_\mathrm{B}}{M_\mathrm{B}m_\mathrm{A}}=K_{\mathrm{f}} b_{\mathrm{B}}$$
    式中，$K_{\mathrm{f}}=\frac{R\left(T_{\mathrm{f}}^{*}\right)^{2}}{\Delta_{\mathrm{fus}} H_{\mathrm{m}}(\mathrm{A})} M_{\mathrm{A}}$，称为凝固点下降常数；$b_\mathrm{B}$为溶液的质量摩尔浓度。
    \begin{table}
        \centering
        \caption{几种溶剂的$K_\mathrm{f}$值}
        \begin{tabular}{cccccc}
            \Xhline{1pt}
            \rowcolor{maincolor!10}
            溶剂 & 水 & 乙酸 & 苯 & 环己烷 & 萘 \\
            \hline
            凝固点$T_{\mathrm{f}}^{*}~/^{\circ}\mathrm{C}$ & 0.00 & 16.60 & 5.53 & 6.5 & 80.25 \\
            $K_\mathrm{f}~/(\si{kg.K.mol^{-1}})$ & 1.86 & 3.90 & 5.12 & 20 & 7.0 \\
            \Xhline{1pt}
        \end{tabular}
    \end{table}
\end{frame}


\begin{frame}{沸点升高}
    沸点(boiling point)是指液体的蒸气压等于外压时的温度。根据拉乌尔定律，在定温时，当溶液中溶质是不挥发性的，溶液的蒸气压总是低于纯溶剂的蒸气压，故溶液的沸点比纯溶剂的要高。
    \\在一定温度下，当溶液中溶剂A与溶液上方的溶剂A的蒸气压达到平衡时，根据相平衡条件，溶剂A在气液两相中的化学势相等
    $$\mu_{\mathrm{A}}^{*}(\mathrm{g})=\mu_{\mathrm{A}}(\mathrm{l})=\mu_{\mathrm{A}}^{*}(\mathrm{l})+R T \ln x_{\mathrm{A}}$$
    若溶液浓度变化$\mathrm{d}X_\mathrm{A}$，沸点相应地变化$\mathrm{d}T$，用与凝固点相同方法处理可得
    $$\Delta T_{\mathrm{b}}=T_{\mathrm{b}}-T_{\mathrm{b}}^{*}=\frac{R\left(T_{\mathrm{b}}^{*}\right)^{2}}{\Delta_{\mathrm{vap}} H_{\mathrm{m}}} \frac{n_{\mathrm{B}}}{n_{\mathrm{A}}}$$
    式中：$T_{\mathrm{b}}^{*}$为纯溶剂的沸点；$T_{\mathrm{b}}$为溶液的沸点；$\Delta_{\mathrm{vap}} H_{\mathrm{m}}$为溶剂的摩尔蒸发焓。
\end{frame}



\begin{frame}{沸点升高}
将溶剂、溶质的量换算为质量摩尔浓度
$$\Delta T_{\mathrm{b}}=\frac{R\left(T_{\mathrm{b}}^{*}\right)^{2}}{\Delta_{\mathrm{vap}} H_{\mathrm{m}}} \frac{M_{\mathrm{A}} m_{\mathrm{B}}}{M_{\mathrm{B}} m_{\mathrm{A}}}=K_{\mathrm{b}} b$$
式中：$K_{\mathrm{b}}=\frac{R\left(T_{\mathrm{b}}^{*}\right)^{2}}{\Delta_{\mathrm{vap}} H_{\mathrm{m}}} M_{\mathrm{A}}$，称为沸点升高常数，它只与溶剂的性质有关。
\begin{table}
    \centering
    \caption{几种溶剂的$K_\mathrm{b}$值}
    \begin{tabular}{ccccccc}
        \Xhline{1pt}
        \rowcolor{maincolor!10}
        溶剂 & 水 & 甲醇 & 乙醇 & 丙酮 & 苯 & 四氯化碳\\
        \hline
        沸点$T_{\mathrm{b}}^{*}~/^{\circ}\mathrm{C}$ & 100 & 64.51 & 78.33 & 56.15 & 80.10 & 76.72\\
        $K_\mathrm{b}~/(\si{kg.K.mol^{-1}})$ & 0.51 & 0.83 & 1.19 & 1.73 & 2.60 & 5.02\\
        \Xhline{1pt}
    \end{tabular}
\end{table}
\end{frame}


\begin{frame}{分配定律}
    定温定压下，在同时共存的两种互不相溶的液体$\alpha$、$\beta$中加入少量某种物质B，当B在$\alpha$、$\beta$两溶剂中溶解达平衡时$\mu_{\mathrm{B}}^{\alpha}=\mu_{\mathrm{B}}^{\beta}$
    $$\mu_{\mathrm{B}}^{\alpha, \ominus}+R T \ln a_{\mathrm{B}}^{\alpha}=\mu_{\mathrm{B}}^{\beta, \ominus}+R T \ln a_{\mathrm{B}}^{\beta}$$
    整理得
    $$\frac{a_{\mathrm{B}}^{\alpha}}{a_{\mathrm{B}}^{\beta}}=\exp [(\mu_{\mathrm{B}}^{\beta \cdot \ominus}-\mu_{\mathrm{B}}^{\alpha, \ominus}) / R T]=K_{a}$$
    溶液很稀时，可用$c_{\mathrm{B}}^{\alpha}$和$c_{\mathrm{B}}^{\beta}$代替$a_{\mathrm{B}}^{\alpha}$和$a_{\mathrm{B}}^{\beta}$，则上式写为$\frac{c_{\mathrm{B}}^{\alpha}}{c_{\mathrm{B}}^{\beta}}=K$(分配定律)，说明在定温定压下溶质在两种共存且不互溶的溶剂中的浓度之比为常数$K$。$K$称为分配常数，是与溶质、溶剂以及温度、压力有关的函数，但压力对其影响较小，可忽略。仅适用于在两溶剂中有相同形态的溶质。    
\end{frame}


\begin{frame}
    \frametitle{分配定律}
    \begin{table}
        \centering
        \caption{\ch{I2}在\ch{H2O}和\ch{CS2}中的分配}
        \begin{tabular}{ccc}
            \bottomrule
            \rowcolor{maincolor!10}
            $c_\mathrm{I_2}(\ch{CS2})$/(\si{g.dm^{-3}}) & $c_\mathrm{I_2}(\ch{H2O})$/(\si{g.dm^{-3}}) & $K=c_\mathrm{I_2}(\ch{CS2})/c_\mathrm{I_2}(\ch{H2O})$ \\ \hline
            0.076 & 0.0017 & 410 \\ \hline
            4.1 & 0.010 & 410 \\ \hline
            6.6 & 0.016 & 410 \\ \hline
            12.9 & 0.032 & 400 \\ \toprule
        \end{tabular}
    \end{table}
\end{frame}


\begin{frame}{萃取}
    用一种与溶液不相混溶的溶剂，从溶液中分离出某种溶质的操作称为萃取(extraction)，所加溶剂为萃取剂。使用分配定律可计算萃取效率。
    \\设用体积为$V_2$的某萃取剂，从体积为$V_1$的含溶质为$m$g的溶液中萃取该溶质，并设溶质分子在萃取剂中与原溶液中形态相同。当分配达平衡时，原溶液中溶质的质量剩$m_1$g，则根据分配定律$K=\frac{\left(m-m_{1}\right) / V_{2}}{m_{1} / V_{1}}$，整理得
    $$m_{1}=m\frac{V_{1}}{K V_{2}+V_{1}}$$
    若用同样体积的萃取剂进行第二次萃取，留在原溶液中溶质的质量$m_2$为
    $$m_{2}=m\left(\frac{V_{1}}{K V_{2}+V_{1}}\right)^{2}$$
\end{frame}


\begin{frame}{萃取}
    经$n$次萃取后，留在原溶液中的溶质为
    $$m_{n}=m\left(\frac{V_{1}}{K V_{2}+V_{1}}\right)^{n}$$
    因为$\frac{V_{1}}{K V_{2}+V_{1}}<1$，所以$n$次萃取后，留在原溶液中的溶质$m_n$已很少。
    \\萃取操作在生产和科研中用途很广，在天然产物有效成分的提取、精制、成分分析等中均有应用。
\end{frame}